\section{}\label{sec:two}
\subsection*{Нестационарная теория возмущений}

\hspace{1em} До этого момента основная теория, которую мы рассматривали, была построена на изменении квантового состояния под действием постоянного гамильтониана. Для описания такой динамики мы решали уравнение Шрёдингера \ref{sec:two}:
\[
i\hbar\frac{d\ket{\psi(t)}}{dt} = \hat{H}\ket{\psi(t)}.
\]
Решение этого уравнения в общем случае записывается как
\[
\ket{\psi(t)} = e^{-\frac{i}{\hbar}Et}\ket{\psi(0)},
\]
где начальное состояние находится из стационарного уравнения Шрёдингера:
\[
\hat{H}\ket{\psi} = E\ket{\psi}
\]

Так как динамика определяется оператором эволюции, то есть комплексной экспонентой, вероятности остаются постоянными, так как при взятии квадрата модуля остаётся только вещественная часть. Получается, такой формализм не позволяет описывать переходы с одного уровня энергии на другой. Для того чтобы разрешить такие переходы, необходимо рассматривать зависящий от времени гамильтониан. Эта зависимость от времени заложена в зависимость от времени потенциала. Далее мы рассмотрим подходы для описания такого подхода и решения задач, связанных с ним.

\subsubsection*{Представление взаимодействия}
Итак, необходимо придумать подход к описанию задачи вида
\[
H = H_0 + V(t),
\]
где $H_0$ отвечает стационарной задаче 
\[
H_0\ket{\psi} = E\ket{\psi}
\]

Для задач такого типа принято использовать \textit{представление взаимодействия}. В этом представлении изменение состояния не целым гамильтонианом, а только его потенциальной частью. Операторы же, в отличие от представления Шрёдингера, тоже начинают эволюционировать под действием стационарной части $H_0$.

Распишем это представление подробнее. Индексом $I$ будем обозначать состояния и операторы в представлении взаимодействия (Interaction), а индексом $S$ -- в представлении Шрёдингера. Тогда состояние в представлении взаимодействия записывается как
\[
\ket{\psi(t)}_I = e^{\frac{i}{\hbar} H_0 t}\ket{\psi(t)}_S
\]
Легко заметить, что в момент $t=0$ эти представления совпадают. Оператор же будет записываться как
\[
V_I = e^{\frac{i}{\hbar} H_0 t} V_S e^{-\frac{i}{\hbar} H_0 t}
\]

Запишем уравнение Шрёдингера в представлении взаимодействия:
\begin{align*}
    i\hbar\frac{\partial}{\partial t}\ket{\psi(t)}_I &= i\hbar\frac{\partial}{\partial t}\left( e^{\frac{i}{\hbar} H_0 t} \ket{\psi(t)}_S \right) =\\
    &= -H_0  e^{\frac{i}{\hbar} H_0 t}\ket{\psi(t)}_S +  e^{\frac{i}{\hbar} H_0 t}\left( H_0 + V_S \right)\ket{\psi(t)}_S =\\
    &=  e^{\frac{i}{\hbar} H_0 t} V_S  e^{-\frac{i}{\hbar} H_0 t}  e^{\frac{i}{\hbar} H_0 t}\ket{\psi(t)}_S=\\
    &= V_I\ket{\psi(t)}_I
\end{align*}
Так, изменение состояния зависит от потенциала $V_I(t)$. Получим динамику произвольного оператора $A_I$, выраженную $H_0$:
\begin{align*}
    \frac{d A_I}{d t} &= \frac{\partial}{\partial t}\left( e^{\frac{i}{\hbar} H_0 t} A_S e^{-\frac{i}{\hbar} H_0 t} \right) =\\
    &= \frac{i}{h} \left(H_0  e^{\frac{i}{\hbar} H_0 t}A_Se^{-\frac{i}{\hbar} H_0 t} + e^{\frac{i}{\hbar} H_0 t} A_S H_0 e^{-\frac{i}{\hbar} H_0 t}\right) +  e^{\frac{i}{\hbar} H_0 t}\frac{\partial A_S}{\partial t}e^{-\frac{i}{\hbar} H_0 t} =\\
    &=  \frac{i}{\hbar}\left[ H_0, A_I \right]  + \left(\frac{\partial A_S}{\partial t}\right)_I
\end{align*}

Используя этот формализм, найдём уравнения для коэффициентов волновой функции $c_n(t)$. Для этого запишем состояние через время $t$, выразив его через линейную комбинацию стационарных состояний $\ket{\psi_n}$:
\[
\ket{\psi(t)}_S = \sum\limits_{n}c_n(t)e^{-\frac{i}{\hbar}E_n t}\ket{\psi_n}
\]
Важно отметить зависимость коэффициентов от времени. Именно это и отражает зависимость гамильтониана от времени, так как в ином случае вся зависимость от времени была бы только в операторе эволюции $U = \exp({-iE_n t/\hbar}$.

Переведем это состояние в представление взаимодействия:
\[
\ket{\psi(t)}_I = \sum\limits_{n}c_n(t)\ket{\psi_n}
\]
Теперь поработаем с уравнением Шрёдингера в представлении взаимодействия. Домножим его слева на $\bra{\psi_n}$:
\[
    i\hbar\frac{\partial}{\partial t}\braket{\psi_n}{\psi(t)} = \bra{\psi_n}V_I\ket{\psi(t)}_I = \sum\limits_{m}\bra{\psi_n}V_I\ket{\psi_m}\braket{\psi_m}{\psi(t)}_I
\]
Для удобства введём следующее обозначение:
\[
\bra{\psi_n}e^{\frac{i}{\hbar} H_0 t}V_S(t)e^{-\frac{i}{\hbar} H_0 t}\ket{\psi_m} = V_{nm}(t)e^{\frac{i}{\hbar}(E_n - E_m)t} = V_{nm}(t)e^{\frac{i}{\hbar}\omega_{nm} t}
\]
Учитывая, что $c_n(t) = \braket{\psi_n}{\psi(t)}$, можно переписать уравнение Шрёдингера в следующем виде:
\[
i\hbar\frac{d}{dt}c_n(t) = \sum\limits_m V_{nm}e^{\frac{i}{\hbar}\omega_{nm}t}c_m(t)
\]
Если записывать это в матричном виде, получим систему ОДУ:
\[
i\hbar\begin{pmatrix} \dot{c_1} \\ \dot{c_2} \\ \vdots \end{pmatrix} = \begin{pmatrix} V_{11} & V_{12}e^{i\omega_{12}t} & \dots \\ V_{21}e^{i\omega_{21}t} & V_{22} & \dots \\ \vdots & \vdots & \ddots \end{pmatrix} \begin{pmatrix} c_1 \\ c_2 \\ \vdots \end{pmatrix}
\]

\subsubsection*{Нестационарная теория возмущений}
До этого момента мы не использовали никаких приближений, связанных с теорией возмущений. К сожалению, задач, которые можно решить точно с помощью такого подхода немного\footnote{вв}. Поэтому, для решения более общего класса задач, воспользуемся теорией возмущений.

Как и в случае со стационарной теорией, нам необходимо разложить коэффициенты в разложении состояния, то есть
\[
c_n(t) = c^{(0)}_n + c^{(1)}_n + \dots
\]
В целом, используя полученную выше систему ОДУ, уже можно иттеративно вычислять различнеы порядки для нестационарной теории возмущений. Для этого на первом шаге достаточно представить $c^{(0)}_n = \delta_{ni}$, где i -- первоначальный уровень, на котором находится система. Решаем уравнения для $c^{(1)}_n$, затем для $c^{(2)}_n$ и так далее.

Рассмотрим ещё один подход, который часто оказывается более удобным для решения связанных с этой темой задач. Для этого снова вспомним, что в данном разделе мы рассматриваем все системы в представлении взаимодействия. В этом представлении удобно будет ввести следующий оператор эволюции:
\[
\ket{\psi(t)}_{I}=U_{i}(t, t_0)\ket{\psi(t_0)}_I
\]

Используя этот оператор, можно записать урванение Шрёдингера следующим образом:
\[
i\hbar\frac{d}{dt} \ket{\psi(t)} = V_{I}(t) \ket{\psi(t)} \implies i\hbar\frac{d}{dt} U_{I}(t, t_0) = V_{I}(t) U_{I}(t, t_0)
\]
Можно решить это операторное дифференциальное уравнение, если поставить начальные условия $\hat{U}_I(t, t_0)\big|_{t=t_0}=\hat{1}$. Используя его, составим из уравнения Шрёдингера интегральное уравнение:
\[
\hat{U}_I(t, t_0) = \hat{1} - \frac{i}{\hbar}\int\limits_{t_0}^{t}\hat{V}_I(t')\hat{U}_I(t', t_0)dt'
\]
Важное условие, что $t_1 > t_2 > \dots > t_n$. Исследование сходимости этого интеграла часто опускается и решение находится численно.

Это уравнение можно аппроксимировать решением, которое носит название \textit{ряд Дайсона} (Dyson series). Выглядит этот ряд следующим образом:
\begin{align*}
\hat{U}_I(t, t_0) = \hat{1} & - \frac{i}{\hbar}\int\limits_{t_0}^{t}\hat{V}_I(t_1)dt_1 + (-\frac{i}{\hbar})^2 \int\limits_{t_0}^{t} dt_1 \int\limits_{t_0}^{t_1}\hat{V}_I(t_1)\hat{V}_I(t_2)dt_2 + \dots\\
&+ (-\frac{i}{\hbar})^n \int\limits_{t_0}^{t} dt_1\dots \int\limits_{t_0}^{t_{n-1}}\hat{V}_I(t_1)\hat{V}_I(t_2)\dots \hat{V}_I(t_n)dt_n +\dots
\end{align*}

Используя этот ряд, попробуем посчитать разные порядки для коэффициентов $c_n (t)$. Пусть система изначально ($t=0$) находится в состоянии $\ket{i}$. Тогда в момент $t$ система будет находиться в состоянии:
\[
\ket{f(t)}_I = U_I(t, 0)\ket{i} = \sum\limits_n\ket{n}\bra{n}U_I(t, 0)\ket{i}
\]
Во-первых, заметим, что $|\bra{n}U_I(t, 0)\ket{i}|^2$ -- это всё та же вероятность перехода системы с уровня $i$ на уровень $n$. Причём эта вероятность совпадает с вероятностью перехода в представлении Шрёдингера\footnote{Важно отметить, что это выполняется только для перехода из одного собственного состояния гамильтониана в другое собственное состояние гамильтониана. В общем случае это не обязательно верно}.

Во-вторых, здесь можно заметить связь с уравнением на коэффициенты $\ket{\psi(t)}_I = \sum_n c_n(t)\ket{\psi_n}$. Действительно, сопаставив их, получми
\[
c_n(t) = \bra{n}U_{I}(t, t_0)\ket{i}
\]
Тогда, можно сопоставить порядки теории возмущения и ряд Дайсона следующим образом: пусть $c_n^{(1)}(t)$ это первый порядок по $V_I(t)$, $c_n^{(2)}(t)$ это второй порядок по $V_I(t)$ и так далее. Тогда, пуолчим следующее выражения для разных порядков:
\begin{align*}
    c^{(0)}_n (t) &= \delta_{ni},\\
    c^{(1)}_n(t) &= \frac{-i}{\hbar} \int\limits_{t_0}^{t}\bra{n}\hat{V}_I(t_1)\ket{i}dt_1 = \frac{-i}{\hbar} \int\limits_{t_0}^{t}e^{i\omega_{ni}t_1}\hat{V}(t_1)_{ni}dt_1,\\
    c^{(2)}_n (t) &= \left(\frac{-i}{\hbar}\right)^2 \sum_m\int\limits_{t_0}^{t}dt_1 \int\limits_{t_0}^{t_1}e^{i\omega_{nm}t_1}\hat{V}_{nm}(t_1)e^{i\omega_{mi}t_2}\hat{V}_{mi}(t_2)dt_2
\end{align*}
Здесь мы ввели стандартные обозначения $V_{ni} = \bra{n}V\ket{i}$ и $\omega_{ni} = (E_n - E_i)/\hbar$. Используя эти приближения, посчитать вероятность перехода из состояния $i$ в состояние $n$ можно по формуле:
\[
P(i \rightarrow n) = |c_n^{(1)}(t) + c_n^{(2)}(t) + \dots|^2
\]
\subsubsection*{Золотое правило Ферми}
Рассмотри пример с постоянным потенциалом следующего вида:
\[
V(t) = \begin{cases}0,&t<0 \\ V, &t\geq0\end{cases}
\]
Хоть потенциал явно не зависит от времени, он в общем случае может зависеть от наблюдаемых $\hat{x},\, \hat{p},\, \hat{s}$.

Выбрав $t_0 = 0$, посчитаем первый порядок поправок к возмущению:
\begin{align*}
    c^{(0)}_n &= \delta_{ni}\\
    c^{(1)}_n &= \frac{-i}{\hbar}\int\limits_{0}^{t}e^{i\omega_{ni} t'}dt' = \frac{V_{ni}}{E_n-E_i}(1-e^{i\omega_{ni}t})
\end{align*}
Посмотрим на квадрат модуля первой поправки:
\[
|c^{(1)}_n|^2 = \frac{|V_{ni}|^2}{|E_n-E_i|^2}(2-2\cos\omega_{ni}t) = \frac{4|V_{ni}|^2}{|E_n-E_i|^2}\sin^2\left(\omega_{ni}t/2\right)
\]

Проанализируем полученное выражение. Для этого сделаем следующее предположение: все интересующие нас конечные состояния лежат в окрестности начального состояния, поэтому мы можем рассматривать их как непрерывный спектр. Хоть это предположение и кажется не естественным, на деле же часто соотвествует реальным экспериментам. Например, при рассеянии на частицах волна может сохранить свою энергию, однако поменять направление. И количество этих направлений непрерывно. Конечно, это работает не всегда, но это именно тот случай, который нам интересно разобрать, так как он имеет полезное практическое применение.

Исходя из этого предположения, введём новую переменную -- плотность конечных состояний как количество уровней внутри интервала энергий ($E,\, E+dE$). Запишем его как $\rho(E)dE$. Эта перемемнная поможет нам перейти от дискретного спектра к непрерывному. Тогда вероятность перехода можно записать как
\[
    \sum\limits_{n:\,E_n\simeq E_i}|c^{(1)}_{n}|^2 \implies \int|c^{(1)}_n|^2\rho(E_n)dE_n=4\int \frac{|V_{ni}|^2}{|E_n-E_i|^2}\sin^2\left[\frac{(E_n-E_i)t}{2\hbar}\right]\rho(E_n)dE_n
\]

Рассмотрим предел этой функции при $t\rightarrow \infty$. Тогда, используя следующее выражение
\[
\lim\limits_{t\rightarrow\infty} \frac{\sin^{2}\alpha x}{\pi \alpha x^2} = \delta(x)
\]
получим следующий предел:
\[
\lim\limits_{t\rightarrow\infty} \frac{1}{|E_n - E_i|^2} \sin^2\left[\frac{(E_n - E_i)t}{2\hbar}\right] = \frac{\pi t}{2\hbar}\delta(E_n - E_i)
\]
Теперь, можно взять интеграл в указанном выше пределе, предварительно вынеся $|V_{ni}|^2$ за знак интеграла, взяв среднее от этого значения:
\begin{align*}
    \lim\limits_{t\rightarrow\infty}\int|c^{(1)}_n|^2\rho(E_n)dE_n &= 4\overline{|V_{ni}|^2}\int \frac{\pi t}{2\hbar}\delta(E_n - E_i)\rho(E_n)dE_n=\\
    &= \frac{2\pi t}{\hbar}\overline{|V_{ni}|^2} \rho(E_n)\bigg|_{E_n\simeq E_i}
\end{align*}

Заметим, что полученное выражение отражает линейную зависимость вероятности от времени. Введём скорость переходов: 
\[
w_{i\rightarrow n} = \frac{d}{dt}\left(\sum\limits_{n}|c^{(1)}_n|^2\right)
\]
Тогда для большого времени, исходя из полученного выше предела, можно получить формулу, названную \textit{золотым правилом Ферми}:
\[
w_{i\rightarrow n} = \frac{2\pi}{\hbar}\overline{|V_{ni}|^2} \rho(E_n)_{E_n\simeq E_i}
\]
Это очень полезная практическая формула, позволяющая посчитать вероятности для перехода из начального состояния в конечные состояния в первом порядке теории возмущений. Напомню, что мы рассматривали постоянный потенциал, и золотое правило применимо именно к нему. Далее мы рассмотрим задачи, где покажем применение этого правила.
\excersize
\begin{center}
\textit{Уровни энергии невозмущенной системы зависят от параметра $\lambda$ и при некотором значении $\lambda_0$ два уровня энергии пересекаются: $E_1^{(0)}(\lambda_0)=E_2^{(0)}(\lambda_0)$. В начальный момент времени система находится на первом уровне в состоянии $\psi_1^{(0)}$. Определите вероятность найти систему в состоянии $\psi_2^{(0)}$ в момент времени $t$, если на систему накладывается постоянное возмущение $\hat{V}$, причем $V_{11} = V_{22} = 0$, но $V_{12} = V^{\ast}_{21} \neq 0$.}
\end{center}

Решение уравнения Шрёдингера $(\hat{H}_0+\hat{V})\ket{\psi(t)}=E\ket{\psi(t)}$ будем искать в следующем виде:
\[
\ket{\psi(t)} = C_1(t)e^{-\frac{i}{\hbar}E_1 t}\ket{\psi_1^{(0)}}+C_2(t) e^{-\frac{i}{\hbar}E_2 t}\ket{\psi_2^{(0)}}
\]

В момент времени $t=0$ $C_1(0)=1, \; C_2(0)=0$, так как в начале система находилась в состоянии $\psi_1^{(0)}$.
\csquare{darklavender}

\[
\ket{\psi_{ALKT}} = \sum\limits_{m_i\in\{1,\, 0,\, -1\}}Tr\left[ A^{m_1}\varepsilon A^{m_2}\varepsilon\dots \varepsilon A^{m_L} \right]\ket{m_1,\, m_2,\, \dots,\, m_L}
\]